\makeabstract
    {
    The Effect of Anesthetic (Tricaine) on Larval Zebrafish Startle Habituation
    }
    {
    Bokei Thompson\textsuperscript{1}, Stacey Beganny\textsuperscript{1}, Josef G. Trapani\textsuperscript{1}
    }
    {
    \textsuperscript{1} Department of Biology, Amherst College, Amherst, MA
    }
    {
    The mechanism of general anesthesia is incompletely understood. Anesthetic effects are dose-dependent and thought to shut down more complex functions (like memory and learning) before basic functions (like reflexes). This study investigates the effect of tricaine, a common fish anesthetic, on larval zebrafish habituation, or reduced response to repeated startle stimuli. Zebrafish startle can be characterized by an extracellular electric field potential; by examining the probability and shape of this field potential between tricaine-anesthesized fish and control fish, this research asks: does low-dose tricaine selectively impair habituation before the startle reflex? More broadly, this research seeks to characterize the state of brain function in an anesthetized subject.
    }
    {
    Zebrafish larvae were divided into control (n=13) and tricaine-anesthesized (n=10) groups. All fish were transgenically modified to express ChR-2 in hair cells, allowing them to optogenetically startle to 480nm blue light. Larvae were head mounted to a dish in low-melt agarose gel; anesthetized fish were incubated in a typical laboratory dose of 32mg/mL tricaine, while controls received a sham incubation in RODI water. After 10 minutes, two electrodes were placed by the head and tail. Larvae were given 5 baseline blue-light startles with 90-second intervals and a habituating series of 100 startles with 5-second intervals. Electric field potential data were recorded and graphed. A new bulk-analysis tool, developed for this project, was used to examine startle probability and amplitude (in mV) of each field potential.
    }
    {
    Anesthesia does not disrupt habituation and causes greater habituation and stronger startles. Anesthetized fish startled 25.6\% less than control fish post-habituation (p=0.067). Surprisingly, anesthetized fish also had a significantly higher field potential amplitude than controls (p\ensuremath{<}0.0001), as well as more variability in field potential amplitude. Since field potentials are a proxy for muscle depolarization, these data mean that anesthetized fish had stronger, less precise muscle responses to startle stimuli than control fish.
    }
    {
    Complex brain functions may not be shut off cleanly during anesthesia. We originally hypothesized that anesthesia would disable habituation in fish while leaving the startle reflex intact. However, the exaggerated habituation and startle intensity found in this research suggests that modulation of the startle reflex is only partially disabled by anesthesia. It is possible that inhibitory components of the startle circuit were impacted, while other activating forces remained. In human anesthesia, the sensitivity of inhibitory and excitatory neurons to anesthesia may also be variable. The way an abnormal combination of active inhibitory and excitatory neurons could impact arousal and sensory processing during anesthesia remains unclear and warrants further exploration.
    }

    \newpage
    